%  ────────────────────────────────────────
%  کامپایل: XeLaTeX (الزامی)
%  حداقل دو بار کامپایل لازم است (برای فهرست مطالب و ارجاعات)
% ============================================================
\documentclass[12pt,a4paper]{book}

% ════════════════════════════════════════
%  ۱. بسته‌های پایه (قبل از xepersian)
%     ⚠ ترتیب بسته‌ها مهم است!
%     ⚠ xepersian باید آخرین بسته باشد!
% ════════════════════════════════════════

% --- رنگ و گرافیک ---
\usepackage{xcolor}
\usepackage{graphicx}
\usepackage{colortbl}          % ← اصلاح: برای \cellcolor

% --- صفحه‌آرایی ---
\usepackage{geometry}
\usepackage{fancyhdr}
\usepackage{setspace}
\usepackage{pdflscape}         % ← اصلاح: برای محیط landscape

% --- جدول ---
\usepackage{longtable}
\usepackage{booktabs}
\usepackage{multirow}          % ← اضافه: برای \multirow در جداول
\usepackage{multicol}
\usepackage{array}             % ← اضافه: برای ستون‌های سفارشی

% --- TikZ (با تمام کتابخانه‌های لازم) ---
\usepackage{tikz}
\usetikzlibrary{
	positioning,                 % ← اصلاح: برای 'of' syntax
	arrows.meta,                 % ← اصلاح: برای {Stealth}
	calc,                        % ← اضافه: برای محاسبات مختصاتی
	shapes.geometric,            % ← اصلاح: برای ellipse و سایر اشکال
	shapes.symbols,              % ← اضافه
	decorations.pathreplacing,   % ← اضافه
	mindmap,                     % ← اصلاح: برای نمودار mindmap
	backgrounds,                 % ← اضافه
	fit                          % ← اضافه
}
\usepackage{pgf-pie}
% --- جعبه‌های رنگی ---
\usepackage[skins,breakable]{tcolorbox}

% --- لیست‌ها ---
\usepackage{enumitem}
\usepackage{amssymb}           % ← برای $\bigstar$, $\blacksquare$, $\blacktriangle$
\usepackage{bm}                % ← برای $\boldsymbol{}$
% --- متفرقه ---
\usepackage{float}
\usepackage{footnote}
\makesavenoteenv{longtable}
\usepackage{etoolbox}
\usepackage{pdfpages}

% --- عنوان‌بندی ---
\usepackage{titlesec}

% --- لینک‌ها (قبل از xepersian) ---
\usepackage{hyperref}

% ════════════════════════════════════════
%  ۲. fontspec و تنظیمات قلم
%     ⚠ باید قبل از xepersian باشد
%     ⚠ از قلم‌هایی استفاده شده که
%       در اکثر سیستم‌ها موجودند
% ════════════════════════════════════════
\usepackage{fontspec}

% --- بررسی وجود قلم لاتین ---
% اگر TeX Gyre Termes ندارید، از Times New Roman یا Latin Modern استفاده شود
\IfFontExistsTF{TeX Gyre Termes}{%
	\setmainfont{TeX Gyre Termes}%  % اگر نصب است
}{%
	\IfFontExistsTF{Times New Roman}{%
		\setmainfont{Times New Roman}% % جایگزین اول
	}{%
		\setmainfont{Latin Modern Roman}% % جایگزین نهایی (همیشه در TeX هست)
	}%
}

% ════════════════════════════════════════
%  ۳. xepersian — آخرین بسته (الزامی!)
% ════════════════════════════════════════
\usepackage{xepersian}

% --- قلم فارسی ---
% IRLotus باید روی سیستم نصب باشد
% جایگزین‌ها: B Nazanin, Vazirmatn, XB Niloofar
\settextfont[Scale=1.05]{IRLotus}
\setdigitfont{IRLotus}

% --- قلم لاتین در متن فارسی ---
% ⚠ setlatintextfont باید بعد از xepersian باشد
\IfFontExistsTF{TeX Gyre Termes}{%
	\setlatintextfont[Scale=0.95]{TeX Gyre Termes}%
}{%
	\IfFontExistsTF{Times New Roman}{%
		\setlatintextfont[Scale=0.95]{Times New Roman}%
	}{%
		\setlatintextfont[Scale=0.95]{Latin Modern Roman}%
	}%
}

% ════════════════════════════════════════
%  ۴. تنظیمات صفحه
% ════════════════════════════════════════
\geometry{
	top=30mm,
	bottom=30mm,
	right=30mm,
	left=25mm,
	headheight=30pt             % ← اصلاح: از 15pt به 30pt
}
\setlength{\parindent}{0pt}
\setlength{\parskip}{8pt}
\onehalfspacing

% ════════════════════════════════════════
%  ۵. رنگ‌های اصلی سند (پالت ملایم)
% ════════════════════════════════════════
\definecolor{cPrimary}{HTML}{1B2A4A}    % سرمه‌ای ملایم
\definecolor{cAccent}{HTML}{C0392B}     % قرمز — هشدارها
\definecolor{cLight}{HTML}{F5F5F0}      % پس‌زمینه‌ی روشن
\definecolor{cGray}{HTML}{7F8C8D}       % خاکستری متن فرعی
\definecolor{cGold}{HTML}{B7950B}       % طلایی — تأکیدهای مثبت
\definecolor{cWarn}{HTML}{E67E22}       % نارنجی — هشدار متوسط
\definecolor{cGreen}{HTML}{27AE60}      % سبز — موفقیت
\definecolor{cConsolid}{HTML}{6C3483}   % بنفش مات — تثبیت و تداوم

% ════════════════════════════════════════
%  ۶. استایل عنوان‌ها
% ════════════════════════════════════════
% اول: فاصله‌ها
\titlespacing*{\section}{0pt}{12pt}{4pt}
\titlespacing*{\subsection}{0pt}{10pt}{3pt}
\titlespacing*{\paragraph}{0pt}{6pt}{1em}

% دوم: فرمت‌ها (باید بعد از spacing باشد)
\titleformat{\chapter}[display]
{\normalfont\Huge\bfseries\color{cPrimary}}
{\chaptertitlename\ \thechapter}{20pt}{\Huge}
\titleformat{\section}
{\normalfont\Large\bfseries\color{cPrimary}}
{\thesection}{1em}{}
\titleformat{\subsection}
{\normalfont\large\bfseries\color{cPrimary!80}}
{\thesubsection}{1em}{}
\titleformat{\paragraph}[runin]
{\normalfont\normalsize\bfseries\color{cPrimary!70}}
{}{}{}

% ════════════════════════════════════════
%  ۷. جعبه‌های اطلاعاتی
% ════════════════════════════════════════

\newtcolorbox{principleBox}[1][]{%
	enhanced, breakable,
	colback=cLight,
	colframe=cPrimary,
	coltitle=white,
	fonttitle=\bfseries,
	title=#1,
	boxrule=0.5pt,
	arc=2pt,
	left=8pt, right=8pt, top=6pt, bottom=6pt,
	shadow={0.5mm}{-0.5mm}{0mm}{black!15}
}

\newtcolorbox{warningBox}[1][]{%
	enhanced, breakable,
	colback=cAccent!5,
	colframe=cAccent!70,
	coltitle=white,
	fonttitle=\bfseries,
	title=#1,
	boxrule=0.5pt,
	arc=2pt,
	left=8pt, right=8pt, top=6pt, bottom=6pt,
}

\newtcolorbox{noteBox}[1][]{%
	enhanced, breakable,
	colback=cGold!5,
	colframe=cGold!70,
	coltitle=white,
	fonttitle=\bfseries,
	title=#1,
	boxrule=0.5pt,
	arc=2pt,
	left=8pt, right=8pt, top=6pt, bottom=6pt,
}

\newtcolorbox{infoBox}[1][]{%
	enhanced, breakable,
	colback=blue!3,
	colframe=cPrimary!50,
	coltitle=white,
	fonttitle=\bfseries,
	title=#1,
	boxrule=0.5pt,
	arc=2pt,
	left=8pt, right=8pt, top=6pt, bottom=6pt,
}

% ════════════════════════════════════════
%  ۸. سربرگ و پابرگ
% ════════════════════════════════════════
\pagestyle{fancy}
\fancyhf{}
\fancyhead[RO]{\small\color{cGray} دادگاه در سایه — وجدان منفصل مسئول}
\fancyhead[LE]{\small\color{cGray} سند پیشنهادی}
\fancyfoot[C]{\small\color{cGray}\thepage}
\renewcommand{\headrulewidth}{0.3pt}
\renewcommand{\footrulewidth}{0pt}

% ════════════════════════════════════════
%  ۹. تنظیم لینک‌ها
% ════════════════════════════════════════
\hypersetup{
	colorlinks=true,
	linkcolor=cPrimary,
	urlcolor=cPrimary!70,
	citecolor=cGold,
	pdfauthor={گروه تدوین},
	pdftitle={دادگاه در سایه — وجدان منفصل مسئول},
	pdfsubject={پیشنهاد تأسیس نظام شبه‌قضایی اپوزیسیون ایران}
}

% ════════════════════════════════════════
%  ۱۰. فرمان‌های کمکی
% ════════════════════════════════════════

% --- تأکید کلیدواژه ---
\newcommand{\keyword}[1]{\textbf{\textcolor{cPrimary}{#1}}}

% --- متن هشدار ---
\newcommand{\warn}[1]{\textcolor{cAccent}{#1}}

% --- متن انگلیسی (با پشتیبانی bidi) ---
% ⚠ \lr از xepersian می‌آید و جهت را به چپ‌به‌راست تغییر می‌دهد
\newcommand{\en}[1]{\lr{\textit{#1}}}
% ────────────────────────────────────────
% اگر می‌خواهید متن انگلیسی italic نباشد:
% \newcommand{\en}[1]{\lr{#1}}
% ────────────────────────────────────────

% ════════════════════════════════════════
%  ۱۱. تنظیمات اضافی bidi و RTL
% ════════════════════════════════════════

% longtable در محیط RTL گاهی مشکل دارد — این خط کمک می‌کند:
\AtBeginEnvironment{longtable}{\small}

% اطمینان از عملکرد صحیح شماره‌گذاری فارسی در جداول:
\SepMark{.}


% ============================================================
%  شروع سند
% ============================================================